\section{System model}

\subsection{The setup}
\begin{figure}[H]
    \centering
    \resizebox{0.75\textwidth}{!}{%
    \begin{tikzpicture}[
        >=Latex,
        dim/.style={<->, thick},
        labelstyle/.style={font=\small},
        magnet/.style={draw, thick, fill=gray!25},
        sensor/.style={draw, thick, fill=blue!10},
        ball/.style={draw, thick, fill=gray!35}
    ]

    % Coordinate principali
    \coordinate (base) at (0,0);
    \coordinate (ball_bottom) at (0,2.2);
    \coordinate (ball_center) at (0,3.0);
    \coordinate (ball_top) at (0,3.8);
    \coordinate (magnet_face) at (0,5.2);

    % Base sensore: bordo superiore allineato con y=0
    \draw[sensor] (-1.0,-0.30) rectangle (1.0,0);
    \node[labelstyle, below] at (0,-0.30) {sensore, base $x=0$};

    % Pallina
    \draw[ball] (ball_center) circle (0.8);
    \fill (ball_center) circle (1.5pt);

    % Scritte spostate più a destra
    \node[labelstyle, right] at (1.6,3.00) {centro pallina, $x+r$};
    \node[labelstyle, right] at (1.6,2.20) {misura sensore, $x$};
    \node[labelstyle, right] at (1.6,3.80) {sommità pallina, $x+2r$};

    % Elettromagnete
    \draw[magnet] (-1.1,5.2) rectangle (1.1,5.8);
    \node[labelstyle, above] at (0,5.8) {elettromagnete};
    \node[labelstyle, right] at (1.15,5.2) {$l$};

    % Linee tratteggiate di riferimento
    \draw[dashed] (-1.7,0) -- (1.7,0);
    \draw[dashed] (-1.7,2.2) -- (1.7,2.2);
    \draw[dashed] (-1.7,3.0) -- (1.7,3.0);
    \draw[dashed] (-1.7,3.8) -- (1.7,3.8);
    \draw[dashed] (-1.7,5.2) -- (1.7,5.2);

    % Quote dimensionali a sinistra
    \draw[dim] (-1.65,0) -- (-1.65,2.2);
    \node[labelstyle, left] at (-1.65,1.1) {$x$};

    \draw[dim] (-1.25,2.2) -- (-1.25,3.0);
    \node[labelstyle, left] at (-1.25,2.6) {$r$};

    \draw[dim] (-1.25,3.0) -- (-1.25,3.8);
    \node[labelstyle, left] at (-1.25,3.4) {$r$};

    \draw[dim] (-2.0,0) -- (-2.0,5.2);
    \node[labelstyle, left] at (-2.0,2.6) {$l$};

    % Air gap
    \draw[dim] (1.75,3.8) -- (1.75,5.2);
    \node[labelstyle, right, align=left] at (1.75,4.5)
    {$d(x)$\\$=l-x-2r$};

    % Forze allineate sull'asse centrale della pallina
    % Forza magnetica (verso l'alto)
    \draw[->, thick, red!70!black] (0,3.05) -- (0,4.45);
    \node[labelstyle, right] at (0.12,4.00) {$F_m$};

    % Forza peso (verso il basso)
    \draw[->, thick, blue!70!black] (0,2.95) -- (0,1.55);
    \node[labelstyle, right] at (0.12,2.10) {$mg$};

    % Forza d'inerzia (verso il basso)
    \draw[->, thick, black] (0,2.65) -- (0,1.95);
    \node[labelstyle, left] at (-0.12,2.35) {$m\ddot{x}$};

    \end{tikzpicture}%
    }
    \caption{Schematic of the plant}
    \label{fig:maglev_plant_scheme}
\end{figure}


The magnetic levitation setup consists of an electromagnet, a steel ball, a light-based position sensor, a current sensor, a power amplifier and a data acquisition board. The control input is the command voltage sent to the amplifier, while the measured outputs are the ball position and the coil current.

The plant can be described through two coupled subsystems. The electrical subsystem determines how the command voltage affects the coil current. The mechanical subsystem determines how the generated magnetic force affects the vertical motion of the ball.

\subsection{Electrical subsystem}
\begin{figure}[ht]
    \centering
    \begin{circuitikz}
        \draw
        (0,0) to[V, v=$v_c(t)$] (0,4)
        to[R, l=$R_c$] (2.5,4)
        to[L, l=$L$] (5,4)
        to[R, l=$R_s$] (7.5,4)
        to[short, i=$i(t)$] (7.5,0)
        -- (0,0);
    \end{circuitikz}
    \caption{Electrical model of the electromagnet circuit with coil resistance $R_c$, sense resistance $R_s$, and inductance $L$.}
    \label{fig:rl_detailed}
\end{figure}

The electromagnet is modelled as a resistor-inductor circuit. If $v_c(t)$ is the voltage applied to the coil, the current dynamics are
\begin{equation}
    v_c(t) = R i(t) + L \frac{\dd i(t)}{\dd t}.
    \label{eq:electrical_rl}
\end{equation}

When the amplifier gain is included, the coil voltage can be written as
\begin{equation}
    v_c(t) = K_g u(t),
\end{equation}
where $u(t)$ is the command voltage generated by the controller. Therefore,
\begin{equation}
    \dot{i}(t) = -\frac{R}{L}i(t) + \frac{K_g}{L}u(t).
    \label{eq:current_dynamics}
\end{equation}

In the implemented architecture, an inner current controller is designed first. When this loop is closed, the electrical dynamics can also be approximated by a first-order transfer function
\begin{equation}
    H_i(s) = \frac{I(s)}{I_{ref}(s)} = \frac{\omega_i}{s + \omega_i},
    \label{eq:current_closed_loop}
\end{equation}
where $\omega_i$ is the bandwidth of the closed current loop.

\subsection{Mechanical subsystem}

Let $x(t)$ be the air gap between the ball and the electromagnet, $m$ the mass of the ball and $g$ the gravitational acceleration. The vertical dynamics are obtained from Newton's law:
\begin{equation}
    m\ddot{x}(t) = mg - F_m(x,i),
    \label{eq:newton}
\end{equation}
where the magnetic force is modelled as
\begin{equation}
    F_m(x,i) = K_m \frac{i^2(t)}{x^2(t)}.
    \label{eq:magnetic_force}
\end{equation}

Substituting \eqnref{eq:magnetic_force} into \eqnref{eq:newton}, the nonlinear mechanical equation becomes
\begin{equation}
    \ddot{x}(t) = g - \frac{K_m}{m}\frac{i^2(t)}{x^2(t)}.
    \label{eq:mechanical_nonlinear}
\end{equation}

\subsection{Nonlinear state-space model}

The nonlinear model is written by choosing the state vector
\begin{equation}
    \vect{x} = \begin{bmatrix} x_1 & x_2 & x_3 \end{bmatrix}^T
    = \begin{bmatrix} x & \dot{x} & i \end{bmatrix}^T.
\end{equation}

Using the amplifier command $u$ as input, the nonlinear state-space model is
\begin{equation}
\begin{cases}
    \dot{x}_1 = x_2, \\
    \dot{x}_2 = g - \dfrac{K_m}{m}\dfrac{x_3^2}{x_1^2}, \\
    \dot{x}_3 = -\dfrac{R}{L}x_3 + \dfrac{K_g}{L}u.
\end{cases}
\label{eq:nonlinear_ss}
\end{equation}

If the inner current loop is considered already closed, the third equation can be replaced by
\begin{equation}
    \dot{x}_3 = -\omega_i x_3 + \omega_i i_{ref}.
    \label{eq:current_loop_state}
\end{equation}

\subsection{Parameters and data values}

The nominal parameters used in the model are listed in \tabref{tab:parameters}. These values are used as the starting point for simulation and controller design.

\begin{table}[H]
\centering
\caption{Nominal parameters of the magnetic levitation setup.}
\label{tab:parameters}
\begin{tabular}{llll}
\toprule
Symbol & Description & Value & Unit \\
\midrule
$L_c$ & Coil inductance & $412.5\times10^{-3}$ & H \\
$R_c$ & Coil resistance & $10$ & $\Omega$ \\
$N_c$ & Number of coil turns & $2450$ &  \\
$l_c$ & Coil length & $0.0825$ & m \\
$r_c$ & Coil radius & $0.008$ & m \\
$K_m$ & Magnetic force constant & $6.5308\times10^{-5}$ & N m$^2$/A$^2$ \\
$R_s$ & Current sense resistance & $1$ & $\Omega$ \\
$r_b$ & Ball radius & $1.27\times10^{-2}$ & m \\
$M_b$ & Ball mass & $0.068$ & kg \\
$T_b$ & Ball travel & $0.014$ & m \\
$l$ & Total length & $0.0394$ & m \\
$g$ & Gravity & $9.81$ & m/s$^2$ \\
$\mu_0$ & Magnetic permeability & $4\pi\times10^{-7}$ & H/m \\
$K_b$ & Position sensor sensitivity & $2.83\times10^{-3}$ & m/V \\
$K_g$ & Amplifier gain & $3$ & V/V \\
\bottomrule
\end{tabular}
\end{table}

\subsection{Equilibrium point}

The equilibrium point is defined by imposing zero velocity and zero acceleration:
\begin{equation}
    \dot{x}_1 = 0, \qquad \dot{x}_2 = 0, \qquad \dot{x}_3 = 0.
\end{equation}

For a selected equilibrium air gap $x_e$, the corresponding equilibrium current is
\begin{equation}
    i_e = x_e \sqrt{\frac{m g}{K_m}}.
    \label{eq:equilibrium_current}
\end{equation}

The equilibrium command voltage is obtained from the electrical steady-state condition:
\begin{equation}
    u_e = \frac{R i_e}{K_g}.
    \label{eq:equilibrium_voltage}
\end{equation}

In the implemented controllers, it is important to distinguish absolute variables from deviation variables:
\begin{equation}
    \Delta x = x - x_e, \qquad \Delta i = i - i_e, \qquad \Delta u = u - u_e.
\end{equation}

\subsection{Linearized model}

The nonlinear model is linearized around the equilibrium point $\left(x_e,0,i_e,u_e\right)$. Using deviation variables, the linearized model is
\begin{equation}
    \Delta\dot{\vect{x}} = A\Delta\vect{x} + B\Delta u,
    \qquad
    \Delta y = C\Delta\vect{x} + D\Delta u.
    \label{eq:linear_ss_general}
\end{equation}

If the raw electrical dynamics are included, the matrices are
\begin{equation}
A =
\begin{bmatrix}
0 & 1 & 0 \\
\dfrac{2g}{x_e} & 0 & -\dfrac{2g}{i_e} \\
0 & 0 & -\dfrac{R}{L}
\end{bmatrix},
\qquad
B =
\begin{bmatrix}
0 \\
0 \\
\dfrac{K_g}{L}
\end{bmatrix}.
\label{eq:linear_matrices_voltage}
\end{equation}

If the inner current loop is already closed and $\Delta i_{ref}$ is the input of the outer-loop model, the matrices become
\begin{equation}
A =
\begin{bmatrix}
0 & 1 & 0 \\
\dfrac{2g}{x_e} & 0 & -\dfrac{2g}{i_e} \\
0 & 0 & -\omega_i
\end{bmatrix},
\qquad
B =
\begin{bmatrix}
0 \\
0 \\
\omega_i
\end{bmatrix}.
\label{eq:linear_matrices_current_loop}
\end{equation}

The measured outputs are position and current:
\begin{equation}
    C =
    \begin{bmatrix}
        1 & 0 & 0 \\
        0 & 0 & 1
    \end{bmatrix},
    \qquad
    D =
    \begin{bmatrix}
        0 \\
        0
    \end{bmatrix}.
    \label{eq:output_matrix}
\end{equation}

\subsection{Nonlinear model validation}

\subsubsection{Time validation}

The nonlinear model is validated by comparing the simulated response with the measured response of the real setup. The same input signal is applied to both systems and the outputs are compared in terms of static gain, transient shape, overshoot and settling time.

% Example figure placeholder
% \begin{figure}[H]
%     \centering
%     \includegraphics[width=0.85\textwidth]{time_validation_nonlinear.pdf}
%     \caption{Time-domain validation of the nonlinear model.}
%     \label{fig:time_validation_nonlinear}
% \end{figure}

\subsubsection{Frequency validation}

The frequency response is validated by applying sinusoidal inputs at different frequencies and comparing the experimental magnitude and phase with the model prediction.

% \begin{figure}[H]
%     \centering
%     \includegraphics[width=0.85\textwidth]{frequency_validation_nonlinear.pdf}
%     \caption{Frequency-domain validation of the nonlinear model.}
%     \label{fig:frequency_validation_nonlinear}
% \end{figure}

\subsection{Linearized model validation}

The linearized model is expected to be accurate only for small deviations around the selected equilibrium point. For this reason, the validation inputs must be small enough to keep the system inside the linear operating range. The linear model is compared with both the nonlinear simulation and the experimental data.
